\normalsize{ \indent
Para ubicar cada una de las actividades en el tiempo,
se optó por un diagrama de Gantt; con propósitos de ver
las fechas críticas sencillamente, en conjunto con la
existencia de un solo camino transitable. A continuación,
se procede a mostrar el diagrama de Gantt de este documento
(se lo ha dividido en dos partes para favorecer su
visualización):
}
\begin{figure}[!ht]
  \caption{Primera parte de planificación de actividades}
  \centering
  \includegraphics[width=\textwidth]{planTrabajoFinal1}
\end{figure}
\begin{figure}[!ht]
  \caption{Segunda parte de planificación de actividades}
  \centering
  \includegraphics[width=\textwidth]{planTrabajoFinal2}
\end{figure}
\newline
\normalsize{ \indent Consideración:}
\begin{itemize}
  \item Todos los manuales se han estado realizando desde
  el inicio de las actividades hasta el final.
\end{itemize}
